\documentclass[a4paper,12pt]{article}

\usepackage[utf8]{inputenc}
\usepackage[T1]{fontenc}
\usepackage[ngerman]{babel}
\usepackage{geometry}
\usepackage{titlesec}
\usepackage{enumitem}
\usepackage{booktabs}
\usepackage{array}
\usepackage{tabularx}
\usepackage{parskip}
\usepackage{hyperref}
\usepackage{lmodern}
\usepackage{microtype}
\usepackage{amssymb}
\usepackage{listings}
\usepackage[simplified]{pgf-umlcd}

% Lange \texttt{}-Bezeichner brechen über \allowbreak in Tabellen; \sloppy
% lockert die Zeilenumbruch-Strafen, damit nichts über den Rand läuft.
\sloppy

\geometry{
  a4paper,
  left=2.5cm,
  right=2.5cm,
  top=3cm,
  bottom=3cm
}

% Section formatting
\titleformat{\section}[block]{\large\bfseries}{Kapitel \thesection:}{0.5em}{}
\titleformat{\subsection}[block]{\normalsize\bfseries}{\thesubsection}{0.5em}{}
\titleformat{\subsubsection}[block]{\normalsize\bfseries}{\thesubsubsection}{0.5em}{}

% Kapitel beginnen immer auf einer neuen Seite.
\let\oldsection\section
\renewcommand{\section}{\clearpage\oldsection}

\setcounter{secnumdepth}{3}
\setcounter{tocdepth}{3}

% Custom colors / highlighting for placeholder text
\usepackage{xcolor}
\newcommand{\placeholder}[1]{\textcolor{gray}{[#1]}}

\lstset{
  basicstyle=\ttfamily\small,
  keywordstyle=\bfseries,
  showstringspaces=false,
  breaklines=true,
  frame=single,
  numbers=none,
  language=Java
}

% List settings
\renewcommand{\labelitemiii}{$\blacksquare$}

\setlist[itemize,1]{label=\textbullet, leftmargin=1.5em}
\setlist[itemize,2]{label=\textopenbullet, leftmargin=1.5em}
\setlist[itemize,3]{label=$\blacksquare$, leftmargin=1.5em}
\setlist[itemize,4]{label=\textbullet, leftmargin=1.5em}

\begin{document}

% ============================================================
% Title Page
% ============================================================
\begin{titlepage}
    \centering
    \vspace*{4cm}

    {\LARGE\bfseries Programmentwurf Konsolen-Schach\par}

    \vspace{3cm}

    \begin{tabular}{ll}
        \textbf{Name:}           & Max Katzenberger     \\[0.5em]
        \textbf{Matrikelnummer:} & 4008573              \\[1em]
        \textbf{Name:}           & Mathis Koeder        \\[0.5em]
        \textbf{Matrikelnummer:} & 3360189              \\[1em]
        \textbf{Abgabedatum:}    & 31.05.2026
    \end{tabular}
\end{titlepage}

% ============================================================
% Kapitel 1: Einführung
% ============================================================
\section{Einführung}

\subsection*{Übersicht über die Applikation}

Die Anwendung ist ein vollständiges Konsolen-Schachspiel im Hot-Seat-Modus. Zwei Spieler teilen sich
einen Rechner und geben ihre Züge nacheinander in algebraischer Notation ein (\texttt{e2-e4},
\texttt{e7-e8=Q}, \texttt{O-O}). Das Brett wird zwischen den Zügen als Unicode-ASCII-Ausgabe dargestellt.

Implementiert sind sämtliche regulären Schachregeln inklusive Rochade, En passant,
Bauernumwandlung sowie die Endspielerkennung für Schachmatt, Patt, 50-Züge-Regel und dreifache
Stellungswiederholung. Partien lassen sich in einer einfachen, von Hand lesbaren Datei speichern und
später wieder laden (Format: FEN-Stellung plus PGN-light-Zugfolge).

Das Projekt löst kein konkretes Praxisproblem, sondern dient als didaktischer Übungsträger für
Clean Architecture, Domain Driven Design, SOLID, Refactorings und Entwurfsmuster im Sinne des
DHBW-Programmentwurfs.

\subsection*{Wie startet man die Applikation?}

Voraussetzungen: JDK 17 oder neuer und ein Terminal mit Unicode-Unterstützung. Der mitgelieferte
Gradle-Wrapper lädt sich Gradle 9 selbst herunter.

\begin{lstlisting}[language=bash]
cd Chess
./gradlew run
\end{lstlisting}

Im Spiel führt der Befehl \texttt{help} eine Übersicht aller Eingaben auf. Eine typische Sitzung:

\begin{lstlisting}
> new Max Mathis
Neue Partie: 4f3c2a-...
> e2-e4
> e7-e5
> save
> quit
\end{lstlisting}

\subsection*{Wie testet man die Applikation?}

Die JUnit-5-Suite läuft komplett über Gradle. Mockito-Tests benötigen das Flag
\texttt{net.bytebuddy.experimental=true}, das im \texttt{build.gradle.kts} bereits
gesetzt ist.

\begin{lstlisting}[language=bash]
./gradlew test
./gradlew jacocoTestReport
open build/reports/jacoco/test/html/index.html
\end{lstlisting}

% ============================================================
% Kapitel 2: Clean Architecture
% ============================================================
\section{Clean Architecture}

\subsection*{Was ist Clean Architecture?}

Clean Architecture nach Robert C. Martin strukturiert ein System in konzentrische Schichten, deren
Abhängigkeiten stets nach innen gerichtet sind. Im Zentrum dieser Schale liegen die
Domänen-Entitäten und Geschäftsregeln, in den äußeren Schichten befinden sich Anwendungs-Use-Cases und
schließlich Adapter wie UI, Datenbank oder Frameworks. Die zentrale Regel \emph{Dependency
Rule} besagt, dass Code aus einer äußeren Schicht zwar auf innere Schichten zugreifen darf, aber
niemals umgekehrt: die Domain weiß nichts von der Datenbank, die Datenbank weiß nichts vom UI.

Im Projekt ergeben sich vier Schichten:

\begin{itemize}
    \item \textbf{domain}: Entitäten (\texttt{Game}, \texttt{Board}, \texttt{Piece}…), Value
          Objects, Domain-Services, Repository-Interfaces.
    \item \textbf{application}: Use-Cases als Application Services (\texttt{GameService},
          \texttt{MoveService}, \texttt{PersistenceService}), Factories, DTOs.
    \item \textbf{infrastructure}: Implementierungen der Repository-Interfaces sowie
          Serialisierer (\texttt{FileGameRepository}, \texttt{FenSerializer}).
    \item \textbf{presentation}: Konsolen-UI mit \texttt{ConsoleApp}, \texttt{BoardRenderer},
          \texttt{InputParser} und Command-Objekten.
\end{itemize}

Abhängigkeitsrichtung: \texttt{presentation $\to$ application $\to$ domain $\leftarrow$ infrastructure}.

\subsection*{Analyse der Dependency Rule}

\subsubsection*{Positiv-Beispiel: Dependency Rule}

\texttt{de.dhbw.chess.domain.entity.Game} ist Aggregat-Wurzel und liegt damit in der innersten
Schicht. Die Klasse importiert ausschließlich aus \texttt{domain.valueobject}, \texttt{domain.entity}
und \texttt{domain.service}, keinen Import in Richtung \texttt{application},
\texttt{infrastructure} oder \texttt{presentation}.

\begin{tikzpicture}
\begin{class}[text width=8cm]{Game}{0,0}
    \attribute{- id : UUID}
    \attribute{- white, black : Player}
    \attribute{- board : Board}
    \attribute{- history : MoveHistory}
    \attribute{- validator : MoveValidator}
    \attribute{- checkDetector : CheckDetector}
    \attribute{- stateEvaluator : GameStateEvaluator}
    \attribute{- activeColor : PieceColor}
    \attribute{- status : GameStatus}
    \operation{+ makeMove(move : Move) : MoveRecord}
    \operation{+ resign(color : PieceColor)}
\end{class}
\end{tikzpicture}

\textbf{Wer hängt von \texttt{Game} ab?} \texttt{GameService} und \texttt{MoveService} aus der
Application-Schicht, sowie alle Command-Klassen der Presentation-Schicht, also ausschließlich
äußere Schichten. \textbf{Wovon hängt \texttt{Game} ab?} Nur von Klassen seiner eigenen Domain.
Damit erfüllt die Klasse die Dependency Rule.

\subsubsection*{Negativ-Beispiel: Dependency Rule}

Im aktuellen Stand existiert keine Domain- oder Application-Klasse, die nach außen importiert. Als
hypothetisches Negativ-Beispiel: würde \texttt{Game} direkt eine \texttt{FileGameRepository}-Instanz
besitzen und beim \texttt{makeMove} synchron schreiben, würde der Domänenkern an die
Infrastruktur-Schicht gekoppelt; ein Wechsel auf eine Datenbank zöge Änderungen mitten in der
Domain nach sich. \texttt{MoveService} in Application als neuer Service.
\texttt{game.makeMove(...)} bleibt in der Domain isoliert, der Service ruft anschließend
\texttt{repository.save(game)} auf.

\subsection*{Analyse der Schichten}

\subsubsection*{Schicht Domain: \texttt{MoveValidator}}

\texttt{de.dhbw.chess.domain.service.MoveValidator} ist ein Domain-Service. Er prüft für eine
Stellung, ob ein Zug regelkonform ist (Bewegungsmuster der Figur, Rochade-Vorbedingungen, En
passant, Selbst-Schach). Es gibt keinen sinnvollen Identitätsbegriff, daher Service statt Entity.
Die Klasse spricht nur Domain-Klassen an (\texttt{Board}, \texttt{Piece}-Hierarchie,
\texttt{MoveHistory}) und ist damit Teil der innersten Schicht.

\begin{tikzpicture}
\begin{class}[text width=8.5cm]{MoveValidator}{0,0}
    \attribute{- checkDetector : CheckDetector}
    \attribute{- castlingRules : CastlingRules}
    \attribute{- enPassantRules : EnPassantRules}
    \operation{+ validate(board, move, color, history)}
\end{class}
\end{tikzpicture}

\subsubsection*{Schicht Infrastructure: \texttt{FileGameRepository}}

\texttt{de.dhbw.chess.infrastructure.persistence.FileGameRepository} implementiert das
Domain-Interface \texttt{GameRepository} und schreibt jede Partie als Textdatei ins
Verzeichnis \texttt{\textasciitilde/.chess-saves/}. Die Klasse kennt die Domain (sie serialisiert
\texttt{Board} und \texttt{MoveHistory}), die Domain kennt aber nicht die Datei.

\begin{tikzpicture}
\begin{interface}[text width=7cm]{GameRepository}{0,0}
    \operation{+ save(game : Game)}
    \operation{+ findById(id : UUID) : Optional<Game>}
    \operation{+ delete(id : UUID)}
\end{interface}
\begin{class}[text width=7cm]{FileGameRepository}{0,-3.5}
    \implement{GameRepository}
    \attribute{- directory : Path}
\end{class}
\end{tikzpicture}

% ============================================================
% Kapitel 3: SOLID
% ============================================================
\section{SOLID}

\subsection*{Analyse Single-Responsibility-Principle (SRP)}

\subsubsection*{Positiv-Beispiel}

\texttt{CheckDetector} hat genau eine Verantwortung: zu prüfen, ob der König einer Farbe auf
einem gegebenen Brett im Schach steht.

\begin{tikzpicture}
\begin{class}[text width=8cm]{CheckDetector}{0,0}
    \operation{+ isInCheck(board : Board, color : PieceColor) : boolean}
    \operation{+ isAttacked(board : Board, pos : Position, by : PieceColor) : boolean}
\end{class}
\end{tikzpicture}

Es gibt nur einen Grund, die Klasse zu ändern: wenn sich die Definition von \glqq{}im
Schach\grqq{} ändert. Sie kennt weder das Aggregat noch Persistenz und mischt keine Verantwortungen.

\subsubsection*{Negativ-Beispiel}

In den Commits \texttt{39706e1} (\glqq{}Game.makeMove with inline validation\grqq{}) und
\texttt{a4d18f5} (\glqq{}inline check detection in Game\grqq{}) hatte \texttt{Game} mehrere
Aufgaben gleichzeitig: Aggregat-Lebenszyklus, Zugvalidierung \emph{und} Schach-Erkennung. Die
gesamte Regel-Logik hing direkt in \texttt{makeMove}, sodass die Klasse mehr als einen Grund hatte,
sich zu ändern. Jedes neue Regeldetail (Rochade, En passant, Selbst-Schach) hätte mitten in
\texttt{Game} eingeflochten werden müssen.

\begin{lstlisting}[caption={Game.makeMove (Stand \texttt{39706e1}): Validierung inline im Aggregat}]
public MoveRecord makeMove(Move move) {
    Objects.requireNonNull(move, "move");
    if (status.isFinal()) {
        throw new IllegalStateException("Partie ist beendet");
    }
    Piece moving = board.pieceAt(move.from());
    if (moving == null) {
        throw new IllegalArgumentException("Kein Stueck auf " + move.from());
    }
    if (moving.color() != activeColor) {
        throw new IllegalArgumentException("Falsche Farbe am Zug");
    }
    List<Position> targets = moving.possibleMoves(board, move.from());
    if (!targets.contains(move.to())) {
        throw new IllegalArgumentException("Zug nicht erlaubt fuer " + moving.type());
    }
    // ... Brett aendern, Historie schreiben, Farbe wechseln
}
\end{lstlisting}

Lösungsweg: \texttt{MoveValidator} und \texttt{CheckDetector} als eigene Domain-Services
extrahieren. \texttt{Game} delegiert dann nur noch und bleibt damit auf seine Aggregat-Verantwortung
beschränkt.

\newpage

\subsection*{Analyse Open-Closed-Principle (OCP)}

\subsubsection*{Positiv-Beispiel}

Die abstrakte Klasse \texttt{Piece} mit der Methode \texttt{possibleMoves(Board, Position)} und den
Subklassen \texttt{Pawn}, \texttt{Rook}, \texttt{Bishop}, \texttt{Queen}, \texttt{Knight},
\texttt{King}.

\begin{tikzpicture}
\begin{abstractclass}[text width=6cm]{Piece}{0,0}
    \attribute{- type : PieceType}
    \attribute{- color : PieceColor}
    \operation{+ possibleMoves(board, from)}
\end{abstractclass}
\begin{class}[text width=3.5cm]{Pawn}{-4,-4}
    \inherit{Piece}
\end{class}
\begin{class}[text width=3.5cm]{Knight}{0,-4}
    \inherit{Piece}
\end{class}
\begin{class}[text width=3.5cm]{Queen}{4,-4}
    \inherit{Piece}
\end{class}
\end{tikzpicture}

Soll eine neue Figur (z. B. \glqq{}Berolina-Bauer\grqq{}) hinzukommen, schreibt man eine neue
Subklasse, ohne bestehende Klassen zu ändern --- offen für Erweiterung, geschlossen für Modifikation.
Das gleiche Prinzip nutzt das Command-Pattern in der Presentation-Schicht: ein neuer CLI-Befehl
ist eine neue \texttt{Command}-Implementation; der Dispatcher in \texttt{ConsoleApp} bleibt
unangetastet.

\subsubsection*{Negativ-Beispiel}

Im Commit \texttt{38655b6} (\glqq{}Piece base class with type-based switch logic\grqq{})
existierte eine einzelne, finale \texttt{Piece}-Klasse mit einer langen
\texttt{switch(type)}-Anweisung in \texttt{possibleMoves}. Jede neue Figur hätte den
\texttt{switch} aufgebrochen. Der Code-Smell wurde im Refactoring 1 (Commit \texttt{4a5b875})
durch Extraktion in Subklassen (Strategy) aufgelöst.

\begin{lstlisting}[language=Java,caption={Piece.possibleMoves (Stand 89e3843): zentrale switch-Logik},basicstyle=\ttfamily\footnotesize,frame=single]
public List<Position> possibleMoves(Board board, Position from) {
    List<Position> moves = new ArrayList<>();
    switch (type) {
        case PAWN:
            addPawnMoves(board, from, moves);
            break;
        case ROOK:
            addSlidingMoves(board, from, moves,
                new int[][] {{1,0},{-1,0},{0,1},{0,-1}});
            break;
        case BISHOP:
            addSlidingMoves(board, from, moves,
                new int[][] {{1,1},{1,-1},{-1,1},{-1,-1}});
            break;
        case QUEEN:
            addSlidingMoves(board, from, moves, new int[][] {
                {1,0},{-1,0},{0,1},{0,-1},
                {1,1},{1,-1},{-1,1},{-1,-1}});
            break;
        case KNIGHT:
            addStepMoves(board, from, moves, new int[][] {
                {1,2},{2,1},{-1,2},{-2,1},
                {1,-2},{2,-1},{-1,-2},{-2,-1}});
            break;
        case KING:
            addStepMoves(board, from, moves, new int[][] {
                {1,0},{-1,0},{0,1},{0,-1},
                {1,1},{1,-1},{-1,1},{-1,-1}});
            break;
        default:
            throw new IllegalStateException("Unbekannter Figurentyp: " + type);
    }
    return moves;
}
\end{lstlisting}

\bigskip

\subsection*{Analyse Liskov-Substitution- (LSP), Interface-Segregation- (ISP),\\
    Dependency-Inversion-Principle (DIP)}

Hier wird das \textbf{Dependency-Inversion-Principle (DIP)} ausgewählt.

\subsubsection*{Positiv-Beispiel}

\texttt{GameService} hängt vom \emph{Interface} \texttt{GameRepository} ab, nicht von einer
konkreten Implementierung. Das Repository wird im Konstruktor injiziert; ob dahinter
\texttt{InMemoryGameRepository}, \texttt{FileGameRepository} oder im Test ein
\texttt{Mockito.mock(GameRepository.class)} steht, bleibt transparent.

\begin{tikzpicture}
\begin{interface}[text width=5.5cm]{GameRepository}{0,0}
\end{interface}
\begin{class}[text width=5.5cm]{GameService}{-4,-3}
    \attribute{- gameRepository}
\end{class}
\begin{class}[text width=5.5cm]{FileGameRepository}{4,-3}
    \implement{GameRepository}
\end{class}
\begin{class}[text width=5.5cm]{InMemoryGameRepository}{4,-6}
    \implement{GameRepository}
\end{class}
\end{tikzpicture}

Die Application-Schicht hängt nur in Richtung Abstraktion (Interface in der Domain), die
Infrastruktur-Schicht implementiert dieselbe Abstraktion --- klassisches DIP.

\subsubsection*{Negativ-Beispiel}

Würde \texttt{ConsoleApp} direkt \texttt{new FileGameRepository(...)} bauen \emph{und} dann
selbst Datei-Pfade konstruieren und FEN parsen, hinge die Presentation-Schicht an einer konkreten
Klasse statt an einer Abstraktion. Aktuell instanziiert \texttt{ConsoleApp} zwar das konkrete
\texttt{FileGameRepository} (Zusammenstellung in \texttt{main}), reicht es aber als
\texttt{GameRepository}-Referenz an die Services weiter, sodass die übrige Anwendung weiterhin
gegen die Abstraktion programmiert ist.

% ============================================================
% Kapitel 4: Weitere Prinzipien
% ============================================================
\section{Weitere Prinzipien}

\subsection*{Analyse GRASP: Geringe Kopplung}

\subsubsection*{Positiv-Beispiel}

\texttt{GameService} kennt nur das \texttt{GameRepository}-Interface und die
\texttt{GameFactory}. Wechselt die Persistenzstrategie, ändert sich am Service nichts. Insgesamt
verläuft Kommunikation in diesem Projekt fast ausschließlich über schmale Interfaces und Value
Objects.

\subsubsection*{Negativ-Beispiel}

Eine engere Kopplung entsteht in der Presentation-Schicht beim CLI-Befehl \texttt{MoveCmd}:
das Command-Objekt kennt nicht nur den \texttt{MoveService}, sondern greift zusätzlich direkt auf
das Domain-Modell zu (\texttt{game.board()}, \texttt{game.activeColor()}, \texttt{King.class}),
um bei einer Rochade-Eingabe (\texttt{O-O}) das Königsfeld zu finden. Damit hängt eine
Presentation-Klasse an mehreren Domain-Klassen gleichzeitig (\texttt{Game}, \texttt{Board},
\texttt{King}, \texttt{Position}). Die Kopplung ist gewollt, aber nicht minimal --- denkbar wäre,
einen Application-Service-Aufruf wie \texttt{moveService.executeCastling(gameId, kingside)}
einzuführen, sodass die CLI keinen Zugriff mehr auf Domain-Entities benötigt.

\begin{lstlisting}[language=Java,caption={MoveCmd: Presentation-Klasse mit direktem Zugriff auf Domain-Entities},basicstyle=\ttfamily\footnotesize,frame=single]
private Position findKing(Game game, PieceColor color) {
    for (int f = 0; f < 8; f++) {
        for (int r = 0; r < 8; r++) {
            Position p = Position.of(f, r);
            if (game.board().pieceAt(p) instanceof King k && k.color() == color) {
                return p;
            }
        }
    }
    throw new IllegalStateException("Kein Koenig auf dem Brett");
}
\end{lstlisting}

\bigskip

\subsection*{Analyse GRASP: Hohe Kohäsion}

\texttt{FenSerializer} hat ausschließlich eine Aufgabe: das Brett-Layout einer Schachpartie
zwischen Java-Objekten und der Standard-FEN-Stellungsnotation zu konvertieren. Beide Methoden
(\texttt{serializePlacement} und \texttt{parsePlacement}) gehören eng zusammen, teilen den
Wertebereich (FEN-Symbole) und werden im selben Use-Case --- Speichern und Laden --- benötigt.
Entsprechend hoch ist die Kohäsion.

\subsection*{Don't Repeat Yourself (DRY)}

Im früheren Stand (Commit \texttt{a4d18f5}) enthielt \texttt{Game} sowohl die
Validierung eines Zugs als auch eine eigenständige Schach-Erkennung. Beide Stellen prüften, ob ein
Probe-Brett den eigenen König in Schach setzen würde --- die Logik (Brett kopieren, Zug ausführen,
Königsfeld auf Angriffsfreiheit prüfen) war doppelt vorhanden.

\begin{lstlisting}[language=Java,caption={Game (Stand \texttt{a4d18f5}): Schach-Prüfung inline und dupliziert},basicstyle=\ttfamily\footnotesize,frame=single]
public MoveRecord makeMove(Move move) {
    // ... (Eigentums- und Zielfeld-Pruefung) ...
    board.move(move.from(), move.to());
    boolean isCheck = isInCheck(activeColor.opposite());
    if (isInCheck(activeColor)) {
        board.move(move.to(), move.from());
        throw new IllegalArgumentException("Zug laesst eigenen Koenig im Schach");
    }
    // ...
}

private boolean isInCheck(PieceColor color) {
    Position kingPos = board.findKing(color);
    if (kingPos == null) return false;
    for (int f = 0; f < Board.SIZE; f++) {
        for (int r = 0; r < Board.SIZE; r++) {
            Position from = Position.of(f, r);
            Piece p = board.pieceAt(from);
            if (p == null || p.color() == color) continue;
            if (p.possibleMoves(board, from).contains(kingPos)) {
                return true;
            }
        }
    }
    return false;
}
\end{lstlisting}

Im Refactoring 2 (Commits \texttt{82a8048}, \texttt{cede2d0}) wurden \texttt{MoveValidator} und
\texttt{CheckDetector} extrahiert; seitdem existiert die Schach-Prüfung an genau einer Stelle und
beide Aufrufer (Validator beim Probespielen, Aggregat beim Erkennen von \glqq{}Schach!\grqq{})
rufen denselben Service auf.

\begin{lstlisting}[language=Java,caption={CheckDetector (Stand \texttt{cede2d0}): Schach-Prüfung an genau einer Stelle},basicstyle=\ttfamily\footnotesize,frame=single]
public class CheckDetector {

    public boolean isInCheck(Board board, PieceColor color) {
        Position kingPos = board.findKing(color);
        if (kingPos == null) return false;
        return isAttacked(board, kingPos, color.opposite());
    }

    public boolean isAttacked(Board board, Position target, PieceColor attacker) {
        for (int f = 0; f < Board.SIZE; f++) {
            for (int r = 0; r < Board.SIZE; r++) {
                Position from = Position.of(f, r);
                Piece p = board.pieceAt(from);
                if (p == null || p.color() != attacker) continue;
                if (p.possibleMoves(board, from).contains(target)) {
                    return true;
                }
            }
        }
        return false;
    }
}
\end{lstlisting}

% ============================================================
% Kapitel 5: Unit Tests
% ============================================================
\section{Unit Tests}

\subsection*{10 Unit Tests}

\bigskip
\begin{tabularx}{\textwidth}{|l|X|X|}
    \hline
    \textbf{Unit Test} & \textbf{Beschreibung} & \textbf{Klasse\#Methode} \\
    \hline
    UT01 & \texttt{Position} weist negative oder zu große Indizes ab. & \texttt{PositionTest\#}\allowbreak\texttt{rejectsOutOfBounds} \\[1.2em]
    \hline
    UT02 & \texttt{PieceColor.opposite()} ist symmetrisch. & \texttt{PieceColorTest\#}\allowbreak\texttt{oppositeIsSymmetric} \\[1.2em]
    \hline
    UT03 & \texttt{Move.equals/hashCode} ist konsistent. & \texttt{MoveTest\#}\allowbreak\texttt{equalsAndHashCode} \\[1.2em]
    \hline
    UT04 & Springer-Sprung-Muster sind korrekt. & \texttt{PieceMovesTest\#}\allowbreak\texttt{knightLShape} \\[1.2em]
    \hline
    UT05 & Bauern-Doppelschritt nur aus Grundreihe. & \texttt{PieceMovesTest\#}\allowbreak\texttt{pawnDoubleStep} \\[1.2em]
    \hline
    UT06 & Kurze Rochade gelingt unter Standardbedingungen. & \texttt{CastlingTest\#}\allowbreak\texttt{kingsideCastling} \\[1.2em]
    \hline
    UT07 & En passant schlägt den passenden Bauern. & \texttt{EnPassantTest\#}\allowbreak\texttt{capturesAdjacentPawn}\allowbreak\texttt{AfterDoubleStep} \\[1.2em]
    \hline
    UT08 & Promotion ersetzt den Bauern durch die gewählte Figur. & \texttt{PromotionTest\#}\allowbreak\texttt{promotionToQueen}\allowbreak\texttt{ReplacesPawn} \\[1.2em]
    \hline
    UT09 & Schachmatt setzt \texttt{GameStatus} korrekt. & \texttt{EndgameTest\#}\allowbreak\texttt{backRankMate}\allowbreak\texttt{EndsGameWithWhiteWin} \\[1.2em]
    \hline
    UT10 & \texttt{GameService} delegiert Persistenz an das Repository. & \texttt{GameServiceMockTest\#}\allowbreak\texttt{startNewGame}\allowbreak\texttt{PersistsViaRepository} \\[1.2em]
    \hline
\end{tabularx}

Insgesamt umfasst die Suite mehr als 80 Tests; die Tabelle zeigt einen repräsentativen Auszug.

\bigskip

\subsection*{ATRIP: Automatic}

Die Tests laufen ausschließlich über \texttt{./gradlew test}, ohne manuelle Eingaben oder externe
Ressourcen. Dateibasierte Tests (\texttt{FileGameRepositoryRoundTripTest}) verwenden
\texttt{@TempDir}, sodass kein gemeinsam genutztes Dateisystem nötig ist. Das Erfolgs-/Fehlersignal
ist der Exit-Code des Gradle-Aufrufs --- automatisierbar in jeder CI.

\subsection*{ATRIP: Thorough}

\textbf{Positiv-Beispiel:} \texttt{PromotionTest} prüft alle vier Promotion-Optionen (Queen, Rook,
Bishop, Knight) sowie zwei Fehlerfälle: Promotion-Pflicht beim Erreichen der letzten Reihe und
Ablehnung von Promotion außerhalb dieser Reihe. Damit deckt der Test den positiven wie den
negativen Pfad gleichmäßig ab.

\textbf{Negativ-Beispiel:} \texttt{InputParserTest\#rejectsMalformedInput} prüft drei Fehlerformen
in einer Methode (\texttt{e2e4}, leerer String, \texttt{null} an \texttt{Position.fromAlgebraic}).
Sauberer wäre eine Aufteilung in dedizierte Tests pro Eingabeklasse, weil der erste Assert beim
Fehlschlagen die übrigen verdeckt.

\subsection*{ATRIP: Professional}

\textbf{Positiv-Beispiel:} \texttt{EndgameTest\#threefoldRepetitionEndsAsDraw} verwendet eine
Zyklus-Schleife mit Statusabfrage, statt eine konkrete Halbzug-Anzahl hart zu kodieren. Damit
bleibt der Test stabil, wenn sich die Implementierung der Wiederholung minimal ändert (etwa: ab
welchem Halbzug der Zähler greift).

\textbf{Negativ-Beispiel:} Die Hilfsmethode \texttt{p(String)} ist in mehreren Test-Klassen
copy-pasted (\texttt{CastlingTest}, \texttt{EnPassantTest}, \texttt{PromotionTest},
\texttt{EndgameTest}). Professioneller wäre eine gemeinsame Test-Utility, etwa \texttt{static
import} aus einer \texttt{Squares}-Hilfsklasse.

\subsection*{Code Coverage}

\texttt{./gradlew jacocoTestReport} erzeugt einen HTML-Bericht unter
\texttt{build/reports/jacoco/test/html}. Die Domain-Schicht ist mit den vorhandenen Tests gut
abgedeckt: alle Figuren-Subklassen, \texttt{MoveValidator}, \texttt{CheckDetector},
\texttt{GameStateEvaluator} sowie \texttt{Game.makeMove} besitzen umfangreiche Pfad-Abdeckung.
Application-Services werden über \texttt{GameServiceMockTest} und \texttt{PersistenceServiceTest}
mit Mockito getestet. Die Presentation-Schicht ist bewusst nur partiell abgedeckt
(\texttt{InputParserTest}); reine Glue-Code-Klassen wie \texttt{ConsoleApp} oder die einzelnen
Command-Objekte werden manuell beim Spielstart verifiziert.

\subsection*{Fakes und Mocks}

\textbf{Mock 1: \texttt{GameRepository}.} In \texttt{GameServiceMockTest} und
\texttt{PersistenceServiceTest} wird das Repository durch \texttt{Mockito.mock(...)} ersetzt. Damit
können wir prüfen, dass der Service \texttt{save(...)} aufruft (\texttt{verify}) und das
\texttt{findById}-Verhalten passend stubben (\texttt{when(repo.findById(id)).thenReturn(...)}),
ohne tatsächlich Dateien zu schreiben.

\begin{tikzpicture}
\begin{interface}[text width=6cm]{GameRepository}{0,0}
\end{interface}
\begin{class}[text width=6cm]{Mockito-Mock}{0,-3}
    \implement{GameRepository}
\end{class}
\end{tikzpicture}

\textbf{Mock 2: \texttt{GameRepository} im PersistenceService-Test.} Auch \texttt{PersistenceService}
wird gegen einen Mock getestet; das Round-Trip-Szenario findet zusätzlich in einem realen
\texttt{FileGameRepository} unter \texttt{@TempDir} statt, sodass Mock und Fake-Verhalten getrennt
voneinander geprüft werden.

% ============================================================
% Kapitel 6: Domain Driven Design
% ============================================================
\section{Domain Driven Design}

\subsection*{Ubiquitous Language}

\bigskip
\begin{tabularx}{\textwidth}{|l|X|X|}
    \hline
    \textbf{Bezeichnung} & \textbf{Bedeutung} & \textbf{Begründung} \\
    \hline
    \texttt{Game}     & Eine vollständige Schachpartie mit Spielern, Brett, Zughistorie, Status. & Der zentrale Begriff aus der Schach-Domäne; sowohl Regelwerk als auch Spieler reden von \glqq{}der Partie\grqq{}. \\[1.2em]
    \hline
    \texttt{Position} & Algebraisches Feld auf dem Brett (\texttt{a1}--\texttt{h8}). & Schachspieler verwenden diese Notation seit Jahrhunderten. Code und UI sprechen denselben Begriff. \\[1.2em]
    \hline
    \texttt{Move}     & Ein einzelner Halbzug (Quelle, Ziel, ggf.\ Promotion). & Die Domänensprache unterscheidet Halbzug von Vollzug; \texttt{Move} entspricht dem englischen Halbzug-Begriff. \\[1.2em]
    \hline
    \texttt{Check}    & Stellung, in der der König eines Spielers angegriffen ist. & Schach\,/\,Check ist ein Kernkonzept der Regeln; im Code modelliert durch \texttt{CheckDetector}. \\[1.2em]
    \hline
\end{tabularx}

\bigskip

\subsection*{Entities}

\texttt{Game} ist eine Entity --- die Identität wird durch eine \texttt{UUID} stabilisiert, der
Zustand (Brett, Status, ziehende Farbe) ändert sich über die Zeit. Zwei Partien mit identischem
Brett sind nicht dieselbe Partie.

\begin{tikzpicture}
\begin{class}[text width=8cm]{Game}{0,0}
    \attribute{- id : UUID}
    \attribute{- board : Board}
    \attribute{- history : MoveHistory}
    \attribute{- status : GameStatus}
    \operation{+ makeMove(move)}
    \operation{+ resign(color)}
\end{class}
\end{tikzpicture}

\subsection*{Value Objects}

\texttt{Position} ist ein klassisches Value Object: unveränderlich, definiert durch seine Werte
(\texttt{file}, \texttt{rank}), Gleichheit über alle Felder. Zwei \texttt{Position.of(4, 1)} sind
identisch, eine \texttt{Position} kennt keine Identität.

\begin{tikzpicture}
\begin{class}[text width=7cm]{Position}{0,0}
    \attribute{- file : int}
    \attribute{- rank : int}
    \operation{+ toAlgebraic() : String}
    \operation{+ fileDelta(other) : int}
\end{class}
\end{tikzpicture}

\subsection*{Repositories}

\texttt{GameRepository} ist die abstrakte Schnittstelle, über die Application-Services Partien
laden und speichern, ohne Persistenz-Details zu kennen. Implementierungen liegen in der
Infrastruktur-Schicht.

\begin{tikzpicture}
\begin{interface}[text width=7cm]{GameRepository}{0,0}
    \operation{+ save(game)}
    \operation{+ findById(id) : Optional<Game>}
    \operation{+ delete(id)}
\end{interface}
\end{tikzpicture}

\subsection*{Aggregates}

\texttt{Game} ist die Aggregat-Wurzel. Sie schließt \texttt{Board}, \texttt{MoveHistory},
\texttt{Player}, \texttt{GameStatus} und die ziehende Farbe als Konsistenzgrenze ein. Sämtliche
regelrelevanten Mutationen (\texttt{makeMove}, \texttt{resign}) erfolgen über \texttt{Game}, sodass
Invarianten wie \glqq{}Status-Wechsel nur nach gültigem Zug\grqq{} oder \glqq{}aktive Farbe
wechselt nach jedem Zug\grqq{} zentral gewahrt bleiben. Externe Klassen erhalten das Brett
nur als Read-Only-Sicht (Snapshot, Renderer).

% ============================================================
% Kapitel 7: Refactoring
% ============================================================
\section{Refactoring}

\subsection*{Code Smells}

\textbf{Code Smell 1 (Switch Statement, Commit \texttt{38655b6}):} ursprünglich enthielt
\texttt{Piece} ein \texttt{switch(type)} mit Zugregeln aller Figuren in einer Methode. Lösung war
\glqq{}Replace Conditional with Polymorphism\grqq{}: Subklassen pro Figurentyp.

\begin{lstlisting}[language=Java,caption={Code Smell 1 -- switch in Piece (Stand \texttt{89e3843})},basicstyle=\ttfamily\footnotesize,frame=single]
public List<Position> possibleMoves(Board board, Position from) {
    List<Position> moves = new ArrayList<>();
    switch (type) {
        case PAWN:   addPawnMoves(board, from, moves); break;
        case ROOK:   addSlidingMoves(board, from, moves, /* ... */); break;
        case BISHOP: addSlidingMoves(board, from, moves, /* ... */); break;
        case QUEEN:  addSlidingMoves(board, from, moves, /* ... */); break;
        case KNIGHT: addStepMoves(board, from, moves, /* ... */);    break;
        case KING:   addStepMoves(board, from, moves, /* ... */);    break;
        default: throw new IllegalStateException("Unbekannter Figurentyp: " + type);
    }
    return moves;
}
\end{lstlisting}

\textbf{Code Smell 2 (Large Class / Duplicated Logic, Commit \texttt{a4d18f5}):} im früheren
Stand vereinte \texttt{Game} die Aggregat-Verantwortung mit Validierung und Schach-Erkennung; die
Schach-Erkennung war zudem an zwei Stellen dupliziert. Lösung war \glqq{}Extract Class\grqq{} zu
\texttt{MoveValidator} und \texttt{CheckDetector}.

\begin{lstlisting}[language=Java,caption={Code Smell 2: Game.makeMove validiert und prüft Schach inline (Stand \texttt{a4d18f5})},basicstyle=\ttfamily\footnotesize,frame=single]
public MoveRecord makeMove(Move move) {
    // ... pseudo-legale Pruefung ...
    board.move(move.from(), move.to());
    boolean isCheck = isInCheck(activeColor.opposite()); // 1. Aufruf
    if (isInCheck(activeColor)) {                        // 2. Aufruf, gleiche Logik
        board.move(move.to(), move.from());
        throw new IllegalArgumentException("Zug laesst eigenen Koenig im Schach");
    }
    // ...
}
\end{lstlisting}

\subsection*{2 Refactorings}

\textbf{Refactoring 1: Replace Conditional with Polymorphism (Commits \texttt{0bf79b1},
\texttt{4a5b875}).}\\
\textbf{Vorher:} eine \texttt{Piece}-Klasse mit \texttt{switch(type)} (siehe Code Smell~1 oben).

\begin{tikzpicture}
\begin{class}[text width=8cm]{Piece (vorher)}{0,0}
    \attribute{- type : PieceType}
    \attribute{- color : PieceColor}
    \operation{+ possibleMoves(...) : List<Position>}
\end{class}
\end{tikzpicture}

\textbf{Nachher:} abstrakte \texttt{Piece}-Basis mit konkreten Subklassen \texttt{Pawn},
\texttt{Rook}, \texttt{Bishop}, \texttt{Queen}, \texttt{Knight}, \texttt{King} (sowie
\texttt{SlidingPiece} als Zwischen-Schicht für Linien-Figuren). Jede Klasse trägt nur ihre eigene
Zuglogik, das System ist offen für neue Figuren ohne Änderung am Bestand (OCP).

\begin{lstlisting}[language=Java,caption={Refactoring 1 nachher (Stand \texttt{4a5b875}): Pawn als eigene Strategy-Klasse},basicstyle=\ttfamily\footnotesize,frame=single]
public abstract class Piece {
    public abstract List<Position> possibleMoves(Board board, Position from);
    // ... gemeinsamer Zustand: type, color, hasMoved ...
}

public final class Pawn extends Piece {
    public Pawn(PieceColor color) { super(PieceType.PAWN, color); }

    @Override
    public List<Position> possibleMoves(Board board, Position from) {
        List<Position> moves = new ArrayList<>();
        int dir = color().isWhite() ? 1 : -1;
        // ... bauern-spezifische Zuglogik ...
        return moves;
    }
}
\end{lstlisting}

\textbf{Refactoring 2: Extract Class für \texttt{MoveValidator} und \texttt{CheckDetector}
(Commits \texttt{82a8048}, \texttt{cede2d0}).}\\
\textbf{Vorher:} \texttt{Game} mit Inline-Validierung und doppelter Schach-Prüfung (siehe Code
Smell~2 oben).

\begin{tikzpicture}
\begin{class}[text width=8cm]{Game (vorher)}{0,0}
    \operation{+ makeMove(move)}
    \operation{- validateInline(move)}
    \operation{- isInCheckInline(...)}
\end{class}
\end{tikzpicture}

\textbf{Nachher:} \texttt{Game} delegiert die Prüfung an einen \texttt{MoveValidator}, dieser
nutzt einen \texttt{CheckDetector} für die Selbst-Schach-Probe. Die Schach-Prüfung existiert
nur noch einmal (DRY), und das Aggregat orchestriert.

\begin{tikzpicture}
\begin{class}[text width=5.5cm]{Game (nachher)}{0,0}
    \attribute{- validator : MoveValidator}
    \attribute{- checkDetector : CheckDetector}
\end{class}
\begin{class}[text width=5.5cm]{MoveValidator}{6,0}
    \attribute{- checkDetector : CheckDetector}
\end{class}
\end{tikzpicture}

\begin{lstlisting}[language=Java,caption={Refactoring 2 nachher (Stand \texttt{82a8048}): MoveValidator als eigener Domain Service},basicstyle=\ttfamily\footnotesize,frame=single]
public class MoveValidator {

    private final CheckDetector checkDetector;

    public MoveValidator(CheckDetector checkDetector) {
        this.checkDetector = checkDetector;
    }

    public void validate(Board board, Move move, PieceColor activeColor) {
        Piece moving = board.pieceAt(move.from());
        if (moving == null)
            throw new IllegalArgumentException("Kein Stueck auf " + move.from());
        if (moving.color() != activeColor)
            throw new IllegalArgumentException("Falsche Farbe am Zug");
        if (!moving.possibleMoves(board, move.from()).contains(move.to()))
            throw new IllegalArgumentException("Zug nicht erlaubt fuer " + moving.type());
        Board probe = board.copy();
        probe.move(move.from(), move.to());
        if (checkDetector.isInCheck(probe, activeColor))
            throw new IllegalArgumentException("Zug laesst eigenen Koenig im Schach");
    }
}
\end{lstlisting}

% ============================================================
% Kapitel 8: Entwurfsmuster
% ============================================================
\section{Entwurfsmuster}

\subsection*{Entwurfsmuster: Strategy}

Die abstrakte Klasse \texttt{Piece} ist Strategie-Träger; jede konkrete Figurklasse
(\texttt{Pawn}, \texttt{Rook}, \texttt{Bishop}, \texttt{Queen}, \texttt{Knight}, \texttt{King})
implementiert eine eigene Variante der \texttt{possibleMoves}-Strategie. Aus Sicht von
\texttt{MoveValidator} und \texttt{Game} ist die Figur ein \texttt{Piece}, die Auswahl der
korrekten Zuglogik geschieht polymorph.

\begin{tikzpicture}
\begin{abstractclass}[text width=6cm]{Piece}{0,0}
    \operation{+ possibleMoves(board, from)}
\end{abstractclass}
\begin{class}[text width=3cm]{Pawn}{-4,-3.5}
    \inherit{Piece}
\end{class}
\begin{class}[text width=3cm]{Rook}{0,-3.5}
    \inherit{Piece}
\end{class}
\begin{class}[text width=3cm]{Knight}{4,-3.5}
    \inherit{Piece}
\end{class}
\end{tikzpicture}

Vorteil: das System ist offen für neue Figuren, der Validator und das Aggregat müssen für eine
neue Figur nicht angefasst werden.

\subsection*{Entwurfsmuster: Command}

Alle CLI-Eingaben sind als Kommando-Objekte modelliert. Das Interface \texttt{Command} hat eine
einzige Methode, der Dispatcher in \texttt{ConsoleApp} wählt anhand des Eingabe-Schlüsselworts das
passende Objekt aus einer \texttt{Map}. Ein neuer Befehl ist eine neue Klasse plus ein
\texttt{put}; bestehende Klassen bleiben unberührt (OCP).

\begin{tikzpicture}
\begin{interface}[text width=5cm]{Command}{0,0}
    \operation{+ execute(session, args)}
\end{interface}
\begin{class}[text width=3cm]{MoveCmd}{-4,-3.5}
    \implement{Command}
\end{class}
\begin{class}[text width=3cm]{SaveCmd}{0,-3.5}
    \implement{Command}
\end{class}
\begin{class}[text width=3cm]{ResignCmd}{4,-3.5}
    \implement{Command}
\end{class}
\end{tikzpicture}

Erweiterungsmöglichkeit: ein zukünftiges Undo lässt sich umsetzen, indem jedes Command beim
Ausführen seinen Inversions-Zustand auf einem Stack ablegt --- die übrige Anwendung muss dafür
nicht umstrukturiert werden.

\end{document}
